\section{Introduction}

Latent world models compress visual observations and predict how their latent
state changes under action.  This makes planning efficient, but it leaves a
basic question unanswered: when a decision depends on an observation that has
already left the model's input window, does the model still retain it?  Low
next-step prediction error does not answer that question.  A target can be
irrelevant to local dynamics yet decisive for a later goal, and a predictor can
fit every intervening frame while discarding it.

Three common evaluations blur this distinction.  Longer context exposes more
raw observations but is not persistent episode state.  Longer imagination
extrapolates farther from its current anchor but cannot recover evidence absent
from that anchor.  A probe can decode a state without showing that a decision
uses it \citep{r_belinkov,r_ravich}.  Recurrent and state-space models, retrieval, and geometric memories
provide useful mechanisms \citep{r_dreamerv3,r_r2i,r_s4wm,r_memworld}; the
interface alone does not show which old evidence survives or affects behavior.

We introduce a controlled audit with four ordered, fail-closed claims:
requirement, cue-time encoding, retention beyond raw context, and use by a fixed
consumer.  Each stronger branch runs only after its prerequisites pass.  Its
thresholds, endpoints, seed grid, and bootstrap rule are locked before outcomes;
a failed branch stops instead of being resplit or relaxed.  Section~2 gives the
operational tests for each rung.

For each admitted branch, the intervention freezes a released LeWorldModel
(LeWM) \citep{r_lewm} or DINO-WM \citep{r_dinowm} checkpoint and trains only a
parameter-matched causal carrier.  Labels never enter carrier training: every
learned arm optimizes the same frozen-host next-feature objective.  The design
compares no persistent state, GRU, LSTM, diagonal state-space, and fixed-trust
predict--correct carriers.  Reads occur before the decision observation, and
the DINO-WM protocols additionally compare full-episode state with the same
carrier reset to the legal three-frame context.  Any retention claim therefore
concerns the frozen-host--carrier composition, not the released checkpoint in
isolation or its native planner.

Our controlled LeWM study asks whether an earlier heterogeneous-task ranking
reversal is reproduced when cue semantics and age are matched.  It is not: the
90\% interval for the specific age-15 fixed-trust-minus-state-space host
interaction lies inside the registered $\pm5$-percentage-point (pp) equivalence
region, while every learned arm's linear-readout accuracy declines with cue
age.  This changes the interpretation of that contrast but does not identify
which unmatched factor produced the earlier reversal; checkpoint, environment,
and training distribution still differ.

The DINO-WM studies show why the full ladder matters.  On PushT, state-space
and fixed-trust carriers preserve highly decodable internal state, yet little
is linearly readable after the frozen predictor at long age; only state-space
clears both age-15 contrasts for six-way binding.  On PointMaze, all four
learned carriers exceed no state and a distribution-shifting context reset,
and one shared external consumer turns full-state reads into resolved
navigation gains.  The same carrier interface therefore supports strong,
weak, and bottlenecked memory depending on the host--task composition.

We make three contributions:
\begin{itemize}
  \setlength\itemsep{1pt}
  \item an ordered, fail-closed protocol that distinguishes necessity,
        cue-time encoding, persistent retention, and executed use;
  \item a cue-and-age-matched cross-host experiment in which the earlier
        ranking reversal is not reproduced, together with complete age curves
        rather than a single-delay score; and
  \item completed extensions to frozen DINO-WM PushT and PointMaze hosts,
        revealing a linear-readout exposure bottleneck and resolved external
        navigation gains under one shared arm-blind consumer.
\end{itemize}
We do not propose a universally superior carrier, and external-consumer results
are not claims about native world-model planning.

\begin{figure}[t]
\centering
\includegraphics[width=\linewidth]{fig_mem_architecture.pdf}
\caption{Our frozen-host intervention.  A real observation becomes either one
LeWM vector or 196 DINO-WM patch trajectories.  Only the green carrier writes
persistent episode state and a zero-initialized residual; the encoder,
predictor, targets, actions, and next-feature objective are shared across arms.
The lower strip separates the legal pre-decision retention read from executed
use by a declared consumer.}
\label{fig:architecture-v2}
\end{figure}

\section{What Counts as Remembering?}

\paragraph{Four ordered claims.}
Let a cue be visible early in a real trajectory and absent from the final raw
context.  We call the number of later real observations its \emph{cue age}.
The audit separates:
\textbf{requirement}, tested by cue-only counterfactuals with identical
post-cue inputs;
\textbf{encoding}, supported by a training-only readout at the last visible cue
frame;
\textbf{retention}, supported at a pre-decision read after shortcut probes
remain below their locked ceilings; and
\textbf{use}, supported only when a fixed consumer changes an executed outcome.
``Full state minus no state'' measures the total carrier benefit.
For DINO-WM, ``full state minus context reset'' estimates, under the reset
intervention, the contribution of carrier state accumulated before the legal
three-frame predictor window while holding carrier parameters fixed.

These claims form a ladder, not interchangeable metrics.  An unencoded cue
cannot diagnose memory.  A retained cue need not appear in the frozen
predictor's preferred coordinates.  A linearly readable state need not be used
by a native planner.  Conversely, a small absolute read can still cause a
paired change in a fixed consumer; we report both quantities so that relative
sensitivity is not mistaken for competent control.

\paragraph{Relation to world-model memory.}
LeWM and DINO-WM predict from bounded latent histories
\citep{r_lewm,r_dinowm}, while Fast-LeWM accelerates imagination without adding
state updated by later real observations \citep{r_fastlewm}.  End-to-end memory
methods span recurrent or structured dynamics \citep{r_dreamerv3,r_r2i,r_s4wm},
recall inside imagined trajectories \citep{r_onmemory}, geometry-indexed or 3D
state \citep{r_memworld,r_persist}, and hybrid or event memories for long-horizon
world-action models \citep{r_memorywam,r_dimwam}.  These works ask how to build
better memory-enabled systems.  A complementary video benchmark treats camera
motion as an intervention on hidden physical-state persistence
\citep{r_statecore}.  Our object of study is different: whether old semantic
evidence is required, stored, exposed, and used after leaving a frozen latent
control host's input window.  RC-aux similarly shows that local predictive
quality need not track planning reachability \citep{r_rcaux}; here the mismatch
concerns evidence over real-observation time.

\section{Controlled Audit Design}

\begin{figure}[t]
\centering
\includegraphics[width=\linewidth]{fig_mem_tasks.pdf}
\caption{Paired counterfactuals and admission.  A and B share one trajectory but
require different answers; only the orange cue pixels differ, as shown by the
$\Delta$ mask, and the teal pre-decision input is identical.  Green cells give
the direct prerequisite interpretation: the visible cue is decodable, legal
late inputs cannot reveal the answer, and no pixels change outside the declared
mask.  These are admission checks, not retention outcomes.}
\label{fig:tasks-v2}
\end{figure}

\subsection{Frozen hosts and a common carrier contract}

The released image encoder, action encoder, latent predictor, and output
projection are immutable.  Each learned carrier receives the current frozen
latent, the intervening action block, and its previous episode state.  Its
read projection is initialized at zero, so every arm begins at the released
host.  GRU, LSTM, diagonal state-space, and fixed-trust carriers are matched in
trainable capacity within each host; optimizer, schedule, data, objective,
fusion point, and state timing are shared.  Fixed-trust uses a channelwise
predicted state and learned input-independent correction trust.  It is not
privileged in the analysis: the protocol, rather than that carrier, is the
methodological contribution.

LeWM consumes three 192-dimensional vector latents and a 10-dimensional action
block.  DINO-WM exposes 196 visual patch trajectories of width 384 together
with proprioceptive and action channels.  One parameter-shared carrier is
applied independently to every patch trajectory; only the frozen native
predictor mixes patches.  The DINO primary read is a fixed $1\!\times\!1$,
$2\!\times\!2$, and $4\!\times\!4$ spatial pyramid of its final visual output.
The carrier's pre-observation prior is preserved as a secondary diagnostic.

\subsection{Counterfactual tasks and stopping gates}

Figure~\ref{fig:tasks-v2} shows the three audit settings.  The completed LeWM
study overlays the same four-color square on Reacher and PushT, with all 16
color--location combinations exactly balanced.  The registered DINO-WM PushT
protocol uses four-way token and six-way ordered-binding tasks on its native
bank; their cue construction and overlay contract appear in
Appendix~\ref{app:dinowm-carrier}.  The PointMaze protocol repeats each native
base window under four transient mini-map goals.  In the matched LeWM study the
cue moves to yield ages 4, 8, and 15 at a fixed endpoint; the DINO-WM protocols
place it at observations 1--3 and move the read to those same ages.  Within
each task, counterfactual variants share every post-cue pixel, action, simulator
state, and endpoint.

Cue-time balanced accuracy must be at least .750; the matched LeWM and
PointMaze protocols additionally require every class recall to reach .700.
Late-input ceilings are chance plus .05: .300 for four-way tasks and .217 for
six-way binding.  A pixel audit requires zero changes outside the declared cue
mask.  Frozen host hashes are checked around every cell.  Failure of any gate
blocks carrier training.  PointMaze also requires its simulator-controller gate
before carrier training, so a failed oracle stops both retention and execution.
Two earlier attempts to create a broader matched suite stopped at their
declared gates and are retained rather than silently discarded.  Their
deterministic selection rule retained color and exact-balanced location; only
then did we construct fresh banks and lock the carrier grid, endpoint,
interaction, and inference before observing any new carrier outcome.  The
color-only comparison is therefore admission-informed rather than a fresh
confirmatory test of target selection.

\begin{table}[!h]
\centering
\footnotesize
\setlength{\tabcolsep}{3.0pt}
\renewcommand{\arraystretch}{1.12}
\caption{Study design, not outcomes.  ``Included'' denotes a predeclared
component; ``conditional'' runs only after its prerequisite gates.  Reset
restarts the same carrier at the legal short context.  Execution uses an
external arm-blind consumer, never a native planner.}
\label{tab:design-v2}
\begin{tabular}{@{}lcccccc@{}}
\toprule
Study & Frozen & Cue-only & Ages & Arms & Reset & Execution \\
\midrule
LeWM / Reacher color & Included & Included & 4/8/15 & 5 & \NA & \NA \\
\rowcolor{AuditGray}
LeWM / PushT color   & Included & Included & 4/8/15 & 5 & \NA & \NA \\
DINO-WM / PushT token + binding & Included & Included & 4/8/15 & 5 & Included & \NA \\
\rowcolor{AuditGray}
DINO-WM / PointMaze goal & Included & Included & 4/8/15 & 5 & Included & Conditional \\
\bottomrule
\end{tabular}
\end{table}

\subsection{Endpoints and inference}

Readouts are standardized multinomial logistic regressions trained only on the
training bank and evaluated once on held-out episodes.  The matched LeWM study
uses 1,200/480 disjoint episodes per host, five carrier seeds, and a separate
readout at each age.  Its 20,000-draw paired bootstrap jointly resamples carrier
seeds across hosts, independently resamples episodes within host, and stratifies
by the 16 color--location combinations.  The registered interaction is the
PushT-minus-Reacher difference in the age-15 fixed-trust-minus-state-space
contrast.  Equivalence uses a two-one-sided test with a $\pm5$-pp margin, hence
the 90\% interval; the 95\% interval answers the distinct nonzero-effect
question.  Operationally, this treats a cross-host change smaller than five
balanced decisions per hundred as negligible for this one contrast; it is not a
domain-wide indifference margin.

DINO-WM PushT uses 1,200/480 native episodes and five seeds for each task
and carrier; its paired bootstrap resamples matched seeds and class-stratified
held-out episodes.  PointMaze uses 300/120 native base windows, expands
each under all four labels, and cluster-resamples intact four-label sets by
native episode as well as matched carrier seed.  No-state copies are
deterministic alignment artifacts and count as one independent model.  The
protocols preserve absolute accuracy, full-minus-no-state, full-minus-reset,
validation predictions, next-feature error, and per-seed arrays; no
heterogeneous score is pooled across hosts.

\section{Matched Cue and Age Do Not Reproduce the Earlier Ranking Reversal}

The motivating LeWM comparison used different cues and delays in Reacher and
PushT and appeared to reverse the ranking of fixed-trust and state-space
carriers.  To test whether that ranking is reproduced under control, we use one
four-way color target on both hosts, balance cue location as an exact nuisance,
and move the cue to ages 4, 8, and 15 while holding each episode endpoint fixed.
A single age-balanced carrier is trained per host, arm, and seed.  All
admissions are recomputed on fresh banks; metrics from the two stopped screens
select the target but never enter this study's inference.

Figure~\ref{fig:matched-v2} gives the complete curves.  Every learned carrier
has an interval-resolved decline from the shortest to the longest cue age on
both hosts, but the slopes and short-age ordering differ by arm and host.  The
no-state read is invariant by construction.  Thus matching does not make the
systems forget at a common rate; it makes the registered ranking comparison
meaningful while exposing the information lost by a single terminal score.

\begin{figure}[t]
\centering
\includegraphics[width=\linewidth]{fig_mem_matched.pdf}
\caption{Matched LeWM result.  Top: the same four-color task forgets with cue
age on both hosts.  Bottom left: fixed-trust has the same age-15 ordering over
state-space on each host (95\% intervals).  Bottom right: only the cross-host
difference is tested for equivalence; its thick 90\% interval lies inside the
registered $\pm5$-pp region, while the thin whisker shows the 95\% interval.}
\label{fig:matched-v2}
\end{figure}

The quantitative anchor is the registered age-15 cross-host interaction,
$-0.21$ pp.  Its 90\% interval lies inside the specified $\pm5$-pp equivalence
region for this fixed-trust-minus-state-space
estimand under the registered resampling model.  It does not imply that the
hosts, carrier dynamics, or architectures are generally equivalent; the margin
and conclusion apply only to this contrast.  Per-location summaries and the
individual host intervals are retained in the appendix.

The registered ranking interaction from the unmatched study is not reproduced
after matching cue semantics and age.  This does not identify the cause of the
earlier reversal: the released checkpoints still differ in environment,
background, dynamics, and training distribution.  The supported conclusion is
narrower: carrier rankings from unmatched tasks cannot be attributed to
architecture alone, and a single delayed endpoint hides substantial
arm-specific age sensitivity under this training contract.

\section{DINO-WM PushT Separates Internal and Host-Exposed Memory}

We next transfer the intervention from vector-latent LeWM to the original
proprioceptive DINO-WM PushT checkpoint.  Figure~\ref{fig:dinowm-carrier}
compares full state with no state and a legal-context reset as cue age grows.

\begin{figure}[t]
\centering
\includegraphics[width=\linewidth]{fig_mem_dinowm_carrier.pdf}
\caption{Persistent carriers on frozen DINO-WM PushT.  Host-output accuracy
(left) and paired gains over no state/reset (middle/right) separate token recall
from ordered binding.  Token advantages fade with age; for age-15 binding only
state-space clears both controls.  The lower-row ordering reflects host
exposure, not internal capacity: fixed-trust retains a highly decodable prior
that this frozen predictor does not expose.  Whiskers are paired 95\%
intervals.}
\label{fig:dinowm-carrier}
\end{figure}

Token recall shows a transition rather than a permanent winner
(Figure~\ref{fig:dinowm-carrier}).  State-space and fixed-trust are clearly
separated from both controls at short cue ages.  At the longest age they still
improve on no state, but neither remains resolved against a carrier restarted
at the legal context.  GRU and LSTM resolve against neither control.  The
sidecar therefore changes the endpoint read, yet the long-age token result no
longer attributes that change specifically to state preceding the reset window.

Six-way binding makes the exposure boundary sharper.  State-space and
fixed-trust both clear the controls at shorter ages, but at age 15 only
state-space remains positive against both.  Its absolute interval still overlaps chance, so the supported
claim is paired sensitivity, not competent six-way recall.  Fixed-trust's
no-state interval touches zero, and GRU is below reset rather than evidence of
retention.

The carrier-prior diagnostic explains why the lower row of
Figure~\ref{fig:dinowm-carrier} is not an intrinsic-capacity ranking.
Fixed-trust retains more linearly decodable binding information internally,
yet state-space is the arm whose output is better exposed through the frozen
predictor and clears both endpoint controls.  In other words, fixed-trust can
keep the cue linearly decodable internally while state-space maps it more
effectively into coordinates used by this host read.  Local next-feature error
also fails to order semantic recall.  The protocol therefore evaluates a
carrier--host composition, not a universal carrier winner.  This native-bank
study covers one DINO-WM host and is neither a matched numerical comparison
with LeWM nor evidence that a native planner uses persistent state.

\section{Retained PointMaze Goals Drive External Navigation}

PointMaze tests a different DINO-WM environment and MuJoCo execution.
Figure~\ref{fig:pointmaze} connects the registered retention contrasts to a
shared external goal consumer and fixed waypoint controller from the pinned
world-model platform \citep{r_stablewm}.

\begin{figure}[t]
\centering
\includegraphics[width=\linewidth]{fig_mem_pointmaze.pdf}
\caption{PointMaze retention and external use.  (a) Balanced accuracy of the
primary frozen-predictor read; chance and no state are 25\%.  (b,c) Paired
full-state gains over no state and legal-context reset.  (d) Age-15 executed
success after a shared arm-blind consumer selects a goal for the fixed waypoint
controller; stars have positive 95\% lower bounds versus both no state and
random goals.  Intervals use 20,000 matched-seed by native-episode cluster
draws; oracle success lies off axis.}
\label{fig:pointmaze}
\end{figure}

Each native base window is repeated under four cue-only mini-map goals while
post-cue pixels, actions, proprioception, and simulator states remain fixed.
All cue, shortcut, frozen-host, and controller prerequisites pass; their exact
values are reported in the appendix.

All four learned carriers resolve against no state and context reset at every
age (Figure~\ref{fig:pointmaze}), although every curve declines as the cue
recedes.  State-space and fixed-trust remain substantially more durable than
the recurrent cells.  Their reset counterfactual also produces a much larger
feature-distribution shift, so full-minus-reset establishes dependence under
that intervention rather than an unconfounded amount of memory.  Unlike PushT,
however, the PointMaze host-facing output remains strongly readable as well as
the carrier prior.  This contrast motivates an exposure hypothesis: the fixed
spatial goal may align better with DINO-WM's spatial-pyramid read than PushT's
token and relational binding variables.  Checkpoint, environment, cue
semantics, and visual statistics co-vary, so the experiment localizes this
difference to the host--task interface but does not identify its cause.

For use, one shared arm-blind multinomial consumer per seed is fitted to pooled
full-state features and its selected goal is executed by the same waypoint
controller.  Every learned arm resolves over both no state and realized random
goals; the largest gain is fixed-trust's $+53.6$ pp.  Goal
selection and execution preserve the carrier ordering, indicating that the
controller does not erase the representation difference.  An independent
implementation refits all five consumers and reproduces their predictions.
This establishes external, consumer-specific use under one frozen checkpoint,
not an executed-reset comparison or native DINO-WM planning.

\begin{table*}[t]
\centering
\scriptsize
\setlength{\tabcolsep}{2.5pt}
\renewcommand{\arraystretch}{1.12}
\caption{Task-level claim ledger. A check denotes passed requirement, cue-encoding, and shortcut gates; carrier cells list arms with a positive registered age-15 paired lower bound. ``None resolved'' means tested but unresolved; \NA\ means not tested. No pooled score or cross-family ranking is computed.}
\label{tab:claim-ledger-v2}
\begin{tabularx}{\textwidth}{@{}>{\raggedright\arraybackslash}p{0.26\textwidth}>{\centering\arraybackslash}p{0.16\textwidth}*{3}{>{\centering\arraybackslash}X}@{}}
\toprule
Study & Gates & vs. no state & vs. reset & Executed use \\
\midrule
LeWM matched color (Reacher / PushT) & \Pass\ 3/3 each & Age curves (both hosts) & \NA & \NA \\
\addlinespace[1pt]
DINO-WM PushT token & \Pass\ 3/3 & State-space, Fixed-trust & None resolved & \NA \\
DINO-WM PushT binding & \Pass\ 3/3 & State-space & State-space & \NA \\
DINO-WM PointMaze goal & \Pass\ 3/3 + controller & All learned & All learned & All learned \\
\bottomrule
\end{tabularx}
\end{table*}


\section{Cross-Study Analysis: Memory Is an Interface Property}

\paragraph{Age changes both level and ranking.}
Every task shows forgetting, but the curves diagnose different transitions.
On PushT, carriers that clear both controls at short age can retain only a
no-state advantage at long age; on PointMaze, the same interfaces remain
resolved against reset.  The matched LeWM curves likewise change ordering and
slope with age.  These within-task shapes are not cross-bank effect estimates,
but they show why one terminal delay is insufficient: an arm can lead early,
converge later, or lose evidence specifically at the host-exposure boundary.

\paragraph{The carrier and frozen host form two memory stages.}
The DINO protocols read both the carrier's pre-observation prior and the frozen
predictor output.  PushT often preserves a highly decodable prior while losing
the cue after frozen host mixing; PointMaze exposes much more of the retained
goal.  Thus a low endpoint can arise because state was not retained, because
the fusion/predictor did not expose it, or both.  The prior is not a substitute
endpoint---a consumer attached there changes the interface---but it prevents
storage and exposure failures from being conflated.

\paragraph{Local predictive fit does not order semantic memory.}
GRU and LSTM can minimize next-feature error while exposing little semantic
memory, and on PointMaze GRU predicts locally better than carriers with much
stronger goal reads.  This is expected from the objective: cue identity can be
irrelevant to local prediction after the cue disappears.  Next-feature loss is
therefore a shared training contract and health diagnostic, not a memory
metric.

\paragraph{Execution identifies a consumer-specific use link.}
On PointMaze, the ordering of goal selection survives conversion into simulated
navigation success.  For this fixed consumer and controller, retained goal
information is therefore behaviorally consequential rather than probe-only.
The result does not imply that a different consumer, native planner, or
end-to-end system would use the same coordinates.  Together, the studies locate
memory in the composition of target, carrier dynamics, frozen exposure
interface, and consumer---not in an architecture name alone.

\section{What the Audit Can and Cannot Localize}

The audit localizes links in a chain rather than one scalar quantity of memory.
Full-versus-reset estimates, under its specified intervention, the contribution
of state accumulated before the native context rather than extra parameters
acting only on the final three frames.  Prior-versus-output gaps localize a
carrier-to-host exposure bottleneck.  External execution establishes that a
fixed consumer can use the representation to change outcomes.  None proves
that a released native planner contains or uses persistent memory, nor does the
study isolate which host, task, or training factor explains the differences.

The resulting failure signatures are distinct.  When full state and reset both
track no state, the carrier supplies no accessible endpoint benefit.  When full
state beats no state but not reset, the module helps without resolved dependence
on state preceding the legal context.  A readable prior paired with a weak host
output instead identifies exposure as the candidate bottleneck.  Finally, an
execution gain establishes use only for the declared consumer.  These cases
cannot be faithfully ordered by one memory score.

Legacy controls sharpen this boundary.  On Reacher, a 56-observation raw
context makes the cue readable while the trained predictor output stays near
chance; in this host, local next-latent fit therefore does not guarantee
semantic exposure (Appendix~\ref{app:rollout}).  Learned rollout cohorts beat
copy-last and respond to shuffled actions.  That pattern is inconsistent with
those two trivial baselines on these cohorts, but their anchor occurs after the
old cue (Appendix~\ref{app:rollout}).  Two label-free repairs optimize their
auxiliary targets yet fail their registered support criteria
(Appendix~\ref{app:repair}).  These are localization results, not evidence that
one alternative objective would universally repair memory.

\paragraph{Diagnosis selects the intervention.}
The decomposition is actionable because different signatures imply different
engineering targets.  If full state, reset, and no state coincide and the prior
is weak, the missing ingredient lies in the write dynamics or training
incentive; changing the planner cannot recover information that was never
stored.  If the prior is strong but the host output is weak, enlarging or
replacing the carrier attacks the wrong boundary: the fusion and exposure path
must preserve coordinates that a downstream read can access.  If the host
output is readable but execution does not change, the unresolved link is the
consumer rather than storage.  Finally, a disruptive reset can reveal state
dependence without measuring a clean quantity of memory, so it should be read
alongside no-state and executed-use contrasts rather than alone.

\paragraph{Architecture comparisons require matched counterfactuals.}
The earlier LeWM reversal illustrates why carrier rankings should be compared
only after target semantics, cue age, endpoint, nuisance balance, host contract,
and inference unit are declared.  A single terminal delay can reward one decay
profile while hiding another; a pooled score can let an easy target mask a
failed one.  Complete age curves show where rankings change, and equivalence
tests should be attached to one practically defined contrast rather than cited
as global architecture equivalence.  We therefore retain task-level no-state,
context-reset, internal-prior, host-output, and executed-use views.  The result
is more diagnostic than one winner label: it says whether to revise the memory
objective, carrier--host interface, or consumer.

\paragraph{The asymmetries generate falsifiable mechanism hypotheses.}
Class count and visible-cue decodability do not explain the PointMaze advantage:
PushT token recall is also four-way and both cues are perfectly readable while
visible.  The separation opens after frozen host mixing.  The binding result is
even more diagnostic: fixed-trust preserves more cue information in its prior,
whereas state-space exposes the stronger endpoint read.  Together these patterns
suggest that alignment between a carrier's read coordinates and the frozen
fusion path can dominate raw storage capacity.  Two controlled tests would
separate this account from task difficulty.  First, a within-checkpoint task
swap can change the remembered relation while holding cue geometry,
trajectories, and spatial readout fixed.  Second, a label-free low-rank exposure
adapter can be trained after freezing the carrier; improved host output with an
unchanged prior would support an exposure bottleneck, whereas a collapsing
prior would return the diagnosis to storage.  These are prospective tests, not
claims established by the present data.

\section{Limitations}

The study covers two LeWM checkpoints and two DINO-WM checkpoints, but its
targets remain controlled categorical cues rather than naturally occurring
long-horizon goals.  It evaluates one frozen sidecar contract, not the native
memory interfaces of DreamerV3, Recall to Imagine, S4WM, Mem-World, or other
recurrent, retrieval, and geometric families; it therefore cannot rank those
families or establish that the audit transfers unchanged to them.  The matched
LeWM banks do not hold checkpoint training distribution or environment fixed,
and the native DINO-WM banks are not bank-matched to LeWM, so their absolute
scores must not be pooled.  In particular, the PointMaze exposure hypothesis
is not identified because host, checkpoint, cue, and environment differ from
PushT.

Parameter matching does not equal optimization matching, and freezing
strengthens causal attribution while excluding end-to-end co-adaptation.
Linear readouts measure accessible information, not native use.  The PointMaze
consumer is an external goal selector followed by a released controller; it
neither tests an executed reset arm nor modifies the native DINO-WM planner.
Five carrier seeds and finite held-out banks leave some rankings unresolved, so
intervals and seed arrays are retained rather than reduced to winner labels.
The reset itself shifts the feature distribution, especially for state-space
and fixed-trust; full-minus-reset identifies dependence on pre-context state
under that reset, not an unconfounded quantity of memory.  The no-state and
external-execution contrasts do not share this reset caveat.  Ages stop at 15
real observations; interference, longer delays, and end-to-end learned
consumers remain outside the present claim boundary.

\section{Reproducibility and Conclusion}

Every formal protocol is hash-locked before outcomes.  Runners reject wrong
GPUs, missing cells, schema drift, altered frozen weights, and nonfinite metrics;
cells commit atomically.  Official verifiers check completeness and provenance,
and a separate implementation reconstructs statistics from preserved
predictions.  The appendix reports every cell, stopped branch, amendment,
bootstrap unit, and claim boundary.

Memory is an empirical chain, not a consequence of recurrence or local fit.
The earlier LeWM ranking interaction is not reproduced after matching;
DINO-WM shows that linear exposure can bottleneck retained state, while
PointMaze converts full-state reads into external success.  We therefore make
bounded task-level claims, not a pooled memory score or universal architecture
winner.
